\chapter{Likelihood-Based Learning with Latent Variables}

\section{Introduction}

Many generative models introduce hidden or latent variables to capture underlying structure in data. While this increases modeling flexibility, it also complicates learning: likelihood functions now involve integrals or sums over unobserved variables.

This chapter develops maximum likelihood learning in latent-variable models. We introduce marginal likelihood, derive the Expectation--Maximization (EM) algorithm, and show how intractability leads naturally to variational approximations.

\section{Why Latent Variables Are Useful}

Latent variables $z$ represent hidden causes or factors explaining observed data $x$.

\begin{itemize}
    \item in clustering, $z$ indicates cluster membership,
    \item in representation learning, $z$ encodes underlying features,
    \item in generative models, $z$ controls variation in outputs.
\end{itemize}

They allow complex distributions to be expressed as mixtures or compositions of simpler ones.

\section{Latent-Variable Models and Marginalization}

\begin{definition}[Latent-Variable Model]
A latent-variable model specifies a joint distribution:
\[
p_\theta(x,z) = p_\theta(x|z)p(z),
\]
where $z$ is unobserved.
\end{definition}

The marginal distribution over observed data is:
\[
p_\theta(x) = \int p_\theta(x,z)\,dz.
\]

\subsection{Marginal Log-Likelihood}

Given data $\{x_i\}_{i=1}^n$, the objective is:
\[
\ell(\theta) = \sum_{i=1}^n \log p_\theta(x_i)
= \sum_{i=1}^n \log \int p_\theta(x_i, z)\,dz.
\]

This integral often makes optimization difficult.

\section{Maximum Likelihood with Hidden Variables}

Unlike fully observed models, the likelihood involves sums or integrals over latent variables:
\[
\log p_\theta(x) = \log \int p_\theta(x,z)\,dz.
\]

This creates two challenges:
\begin{itemize}
    \item \textbf{inference:} computing the posterior $p_\theta(z|x)$,
    \item \textbf{learning:} maximizing the marginal likelihood.
\end{itemize}

These are coupled problems: learning requires inference, and inference depends on the model.

\section{The EM Algorithm: Intuition}

EM alternates between estimating latent variables and updating parameters.

\begin{itemize}
    \item \textbf{E-step:} estimate distribution over latent variables,
    \item \textbf{M-step:} maximize expected complete-data log-likelihood.
\end{itemize}

\subsection{Key Idea}

Instead of optimizing $\log p_\theta(x)$ directly, EM optimizes a surrogate objective based on the current estimate of latent variables.

\section{Lower-Bound Viewpoint via Jensen}

Let $q(z)$ be any distribution. Then:
\[
\log p_\theta(x)
= \log \int q(z)\frac{p_\theta(x,z)}{q(z)}\,dz
\ge
\int q(z)\log \frac{p_\theta(x,z)}{q(z)}\,dz.
\]

\begin{definition}[Evidence Lower Bound (ELBO)]
\[
\mathcal{L}(q,\theta)
=
\mathbb{E}_{q(z)}[\log p_\theta(x,z)]
-
\mathbb{E}_{q(z)}[\log q(z)].
\]
\end{definition}

This provides a lower bound on $\log p_\theta(x)$.

\subsection{Interpretation}

\begin{itemize}
    \item the bound becomes tight when $q(z)=p_\theta(z|x)$,
    \item EM can be viewed as alternating optimization of this bound.
\end{itemize}

\section{The EM Algorithm: Derivation}

EM alternates:

\subsection{E-step}

Set:
\[
q(z) = p_{\theta^{\text{old}}}(z|x).
\]

\subsection{M-step}

Maximize:
\[
\theta^{\text{new}} = \arg\max_\theta
\mathbb{E}_{z \sim p_{\theta^{\text{old}}}(z|x)}
[\log p_\theta(x,z)].
\]

\subsection{Monotonicity}

Each iteration increases the log-likelihood:
\[
\ell(\theta^{\text{new}}) \ge \ell(\theta^{\text{old}}).
\]

Thus EM is guaranteed to improve (or maintain) likelihood at each step.

\section{Worked Example: Gaussian Mixture Model}

Consider a mixture of $K$ Gaussians:
\[
p_\theta(x) = \sum_{k=1}^K \pi_k \, \mathcal{N}(x \mid \mu_k, \Sigma_k).
\]

Introduce latent variables $z \in \{1,\dots,K\}$.

\subsection{E-step}

Compute responsibilities:
\[
\gamma_{ik}
=
p(z=k \mid x_i)
=
\frac{\pi_k \mathcal{N}(x_i \mid \mu_k, \Sigma_k)}
{\sum_j \pi_j \mathcal{N}(x_i \mid \mu_j, \Sigma_j)}.
\]

\subsection{M-step}

Update parameters:
\[
\mu_k = \frac{\sum_i \gamma_{ik} x_i}{\sum_i \gamma_{ik}},
\qquad
\pi_k = \frac{1}{n} \sum_i \gamma_{ik}.
\]

This alternates soft clustering and parameter updates.

\section{Limits of Exact Inference}

In many models, computing $p_\theta(z|x)$ is intractable:
\begin{itemize}
    \item high-dimensional latent spaces,
    \item complex dependencies,
    \item non-conjugate models.
\end{itemize}

This makes exact EM impractical.

\subsection{Consequences}

\begin{itemize}
    \item need for approximate inference,
    \item motivates variational inference,
    \item leads to modern deep generative models.
\end{itemize}

\section{Why Latent Variables Complicate Optimization}

Latent variables introduce:
\begin{itemize}
    \item non-convex objectives,
    \item multiple modes,
    \item dependence on inference quality.
\end{itemize}

Optimization becomes a joint problem of:
\begin{itemize}
    \item fitting parameters,
    \item inferring hidden structure.
\end{itemize}

Thus learning is more complex than in fully observed models.

\section{A Unifying Perspective}

Latent-variable models express complex distributions through hidden structure. Maximum likelihood requires marginalizing over latent variables, which leads to intractability.

EM provides a principled approach by alternating between inference and optimization, while the ELBO offers a variational perspective. These ideas form the foundation for modern methods such as variational autoencoders.

\section{Summary}

In this chapter we developed likelihood-based learning with latent variables. We defined latent-variable models and marginal likelihood, and explained why hidden variables complicate optimization. We derived the EM algorithm and interpreted it through a lower-bound (ELBO) perspective.

We also worked through Gaussian mixtures and discussed the limits of exact inference. These ideas provide the conceptual bridge to variational inference and modern generative models.